% Інтерпретація завершених результатів без нових вимірювань чи міток.
\section{Обговорення}\label{sec:discussion}

\paragraph{Морфологія (RQ1).}
Виправлене розширення лемами підвищує nDCG@10 BM25 на 0,031 за основними мітками з парним 95\% інтервалом від 0,013 до 0,051. Це підтримує приріст для цього набору запитів за цих оцінок. Слабші свідчення за Gemma~4 12B обмежують висновок, незалежний від судді. Попереднє погіршення до виправлення аналізатора також ілюструє залежність від реалізації: розширення неправильно вгаданою лемою може змінювати пошук у протилежний бік. Результат стосується перевіреного правила аналізатора й корпусу, а не морфології як універсально корисної операції.

\paragraph{Злиття й технічні запити (RQ2--RQ3).}
Адаптивне злиття перевищує BM25 за основною метрикою, проте виявленої переваги над EG2 чи RRF немає. E5 має найвище середнє nDCG@10, зокрема в контролі на однакових текстових кандидатах. Відсутність значущого після поправки Холма контрасту E5--adaptive не встановлює однакової ефективності. Абляція переранжування фраз чисельно тотожна повній системі на заморожених прогонах; абляції ідентифікаторів і коротких фрагментів також не підтверджують поліпшення на цій вибірці. Правило може виправляти правдоподібний збій, який трапляється тут недостатньо часто для поліпшення сукупної метрики. Загальні рекомендації щодо правила злиття тому потребували б незалежно написаних запитів і ширшого охоплення цілей.

\paragraph{Зміна охоплення цілей.}
Окреме розширення дає інший порядок: EG2 та адаптивне злиття мають вищі точкові оцінки за E5 за його початковими мітками Gemini. Між наборами також змінюються автор запитів, формати цілей і варіант оцінювання. Тому розширення не визначає, який чинник спричиняє цю відмінність, а його середні не слід вважати поліпшенням «до/після» відносно початкового тесту. Збереження аналізу всіх 50 записів поряд із детермінованою перевіркою на 49 формулюваннях утримує початкові дані й робить повтор формулювання явним.

\paragraph{Мультимодальне зіставлення (RQ4).}
Завдання «підпис--зображення» та «транскрипт--запис» перевіряють міжмодальне зіставлення двома мовами за парами з джерельних наборів. Повнота на першій позиції становить 0,857 для українських підписів і 0,889 для російських серед 1\,000 зображень. Це окремі результати від текстового пошуку з модельними оцінками. Для мовлення використання того самого транскрипту як запиту й індексованого тексту є умовою верхньої межі; навіть міжмовні транскрипти --- паралельні джерельні речення. Жодне завдання не встановлює ефективності для людського перефразування, довгого запису чи довільних фотографій. Вони підтверджують реалізований зв'язок модальностей у зазначених умовах зіставлення.

\paragraph{Квантування та охоплення оцінками (RQ5).}
Звичайне порівняння Q4 має від'ємну точкову різницю, тоді як конденсоване оцінювання не показує такого самого погіршення; 9,9\% результатів Q4 у першій десятці не оцінено. Тому звичайна метрика за замороженим пулом не ідентифікує внутрішньої втрати релевантності Q4. Водночас малі відмінності f16/Q8 та інтервали, що охоплюють нуль, не встановлюють еквівалентності. Бекенд і точність змінюються разом, а для 21 рядка PDF-сторінок змінюється ще й вхідна модальність. Якість, розмір файлів і витрати кодування слід читати разом із цими умовами охоплення та реалізації, а не як ізольований причинний ефект квантування.

\paragraph{Витрати фактичного пошукового шляху.}
На двох мільйонах синтетичних векторів медіана фактичного шляху ранжування архіву становить близько 2\,442~мс, а оптимізованого еталона NumPy для першої десятки --- близько 107~мс. Шляхи виконують різну роботу: метод архіву збирає кандидатів, повністю сортує оцінки й формує 400 кортежів. Різниця виявляє роботу реалізації, яку приховав би опис лише множення матриці на вектор. HNSW із глибиною пошуку 256 витрачає близько 1,28~мс і має recall@10 за точним збігом ідентифікаторів 0,9825 на цій синтетичній колекції. Послідовні вимірювання у звичайній робочій сесії не є оцінкою наскрізного чи парного прискорення; повнота може не переноситися на реальний архів такого розміру. Це обґрунтовує приблизну індексацію та відбір кандидатів як подальші напрями інженерної роботи.

\paragraph{Оцінювання та практична інтерпретація.}
Узгодженість суддів, перевірки цільових документів і навмисно безвідповідні запити характеризують поведінку оцінювання, проте не встановлюють відповідності людським судженням. Перегляд Gemini зберігає залежність від початкових суддів і видимих оцінок, тому точкові результати супроводжуються аналізом чутливості до міток. Незалежне людське дослідження та запити людей, які шукають власні документи, залишаються майбутньою перевіркою й не подаються тут як виконані. Перевірки приватного архіву показують задоволення операційних очікувань автора в цьому середовищі, проте повторне використання під час розроблення не дозволяє заявляти ефективність для користувачів. Разом ці свідчення дають доступний для перевірки опис якості й витрат на одному побудованому бенчмарку та в одному робочому архіві; найсильніші порівняння залежать від міток, вибірок і меж вимірювання.
